\subsection{Liver eQTL colocalization defines a 430-gene multi-signal map across 35 liver-trait GWAS strata}\label{liver-eqtl-colocalization-defines-a-provisional-436-gene-map-across-35-liver-trait-gwas-strata}

An earlier variant-identity check identified duplicate position identifiers in the saved eQTL fits of 90 of 436 r2 primary genes, alongside negative eigenvalues (down to $-$0.62) in pairwise-missing 1000 Genomes LD matrices. Allele-aware, positive-semidefinite (PSD) LD-projected refitting of all 90 genes across the 35 primary GWAS strata (\texttt{EQTL-DEDUP-PSD-20261118}; \texttt{max\_iter = 5000}) passed parity on 5 of 5 unaffected control genes (PNPLA6, MED17, CEP68, GCLC, PKM; max \textbar $\Delta$PP.H4\textbar{} = 2.55 $\times$ 10$^{-4}$) and retained 84 of 90 genes as positive (recovering IRF5 [0.955] and ABHD12 [0.714] from non-convergence, removing 4 lost genes [CDK5RAP3, KANSL1-AS1, LRRC37A, THBS3], and leaving 2 untestable [CTSS and the unresolved repeat-length source event at MPRIP]). Unless explicitly labeled as historical r2 denominators, all primary genetic counts, figures and Unified Gene Catalog entries below reflect this 430-gene multi-signal set and 990-gene union (Methods; Supplementary Note 1).

To map inherited variation relevant to MASLD, we assembled a prespecified set of direct-disease, MRI-derived liver-fat, and liver-enzyme GWAS from several ancestry groups (Figure 2A) (Namjou et al., 2019; Anstee et al., 2020; Ghodsian et al., 2021; Liu et al., 2021; Haas et al., 2021; Sinnott-Armstrong et al., 2021; Kanai et al., 2018; Sveinbjornsson et al., 2022; van der Meer et al., 2022; Kurki et al., 2023; Verma et al., 2024; Karczewski et al., 2025). We kept direct disease and imaging traits apart from the ALT, AST, and GGT proxy phenotypes. Enzyme associations report liver biology but are not equivalent to direct MASLD susceptibility. The registry contains 35 tier-1/2 strata from five GWAS ancestries and six liver traits. Fine-mapping retained 360 trait-specific loci, which correspond to 265 globally unique loci (Wang et al., 2020).

Colocalization against the Broadaway liver-eQTL panel (Broadaway et al., 2024) gave two maps that differ in what they can observe (Figure 2B). As in Figure 1, each result records phenotype applicability, eQTL and LD coverage, method eligibility, evidence state, unresolved alternative, and next discriminating experiment. Both methods report PP.H4, the posterior probability that the GWAS and eQTL associations share one variant. Single-variant ABF and multi-signal SuSiE colocalization have different eligibility requirements (Giambartolomei et al., 2014; Wallace, 2021). SuSiE needs converged eQTL and GWAS fits and usable LD, so a missing multi-signal result marks limited observability, not a negative result.

\begin{mainfigureplate}{2}{650bp}
\noindent\makebox[\linewidth][c]{\paperpanel{A}{340bp}{figures/main/figure2/Fig2A_gwas_alluvial.pdf}\hspace{10bp}\paperpanel{B}{170bp}{figures/main/figure2/Fig2B_PIP_vs_PPH4_clean.pdf}\hspace{10bp}\paperpanel{C}{110bp}{figures/main/figure2/Fig2C_lead_class_clean.pdf}}
\par\vspace{8pt}
\noindent\makebox[\linewidth][c]{\paperpanel{D}{180bp}{figures/main/figure2/Fig2D_ancestry_treemap.pdf}\hspace{10bp}\paperpanel{E}{198bp}{figures/main/figure2/Fig2F_crossancestry_coloc.pdf}\hspace{10bp}\begin{minipage}[t]{240bp}\vspace{0pt}\paperpanel{F}{200bp}{figures/main/figure2/Fig2F_candidate_trait_support.pdf}\par\vspace{10pt}
\paperpanel{G}{210bp}{figures/main/figure2/Fig2G_candidate_followup_modalities.pdf}
\end{minipage}}

\end{mainfigureplate}
\papercaption{2}


After the PSD eQTL refit, the genetic union contains 990 genes: 430 primary multi-signal genes (PP.H4 bands \textgreater{} 0.9 / 0.8--0.9 / 0.5--0.8: 239 / 89 / 102) and 560 single-signal-only genes (67 / 89 / 404; Figure 2A; Methods). Three genes drawn as single-signal only (CKB, GSTM2 and LILRB3) are SuSiE-untestable and are not SuSiE negatives. Because the two methods can test different genes, we report each count against the genes that method could test, and we report single-signal-only genes separately. We ran colocalization twice. The second run repeated only the positives from the first run and added the shared-posterior check, which requires both signals in a pair to keep at least half of their posterior probability on variants that the GWAS and eQTL fits share. Negative results from the first run were carried forward unchanged. The count of finite results therefore includes these carried negatives, and we report it separately from the subset that passed the shared-posterior check. The GWAS fit is built on the variants of the eQTL fit, so only the eQTL signal can fail this check. As a sensitivity analysis, we applied the check to the 12,087 carried-forward negative gene-by-GWAS tests in the tier-1/2 studies. In 343 of them (2.8\%), every signal pair failed, so the test became untestable; in 1,744, only some signal pairs failed and the call stayed the same; the other 10,000 passed (Methods). A carried negative that fails the check becomes untestable, not positive, so the colocalized genes do not change. Of the 436 r2 primary SuSiE genes, 424 were protein-coding and 12 were lncRNA targets (after the PSD refit, 419 are protein-coding and 11 are lncRNA). Gene-level eQTL data do not identify transcript isoforms or give a biotype-specific enrichment denominator.

Fourteen primary genes were selected for diagnostics because their only support came from 18 near-one SuSiE fits. Fourteen of those fits have an LD-inconsistent driving variant, and removal of diagnostic outliers eliminates support for eleven genes. This post hoc sensitivity can also remove genuine complex signals, so it does not prove false colocalization. In addition, individual gene-level colocalizations such as RORA (BBJ GGT) and FABP1 (ALT) remain subject to local LD hitchhiking with neighboring signals and are interpreted as locus-anchored candidates rather than isolated causal genes (Supplementary Note 1).

The COLOC gene counts are not derived from the 360 fine-mapped loci. Of the 436 r2 multi-signal genes, 100 had a colocalized lead variant within 250 kb of a fine-mapped window from the same GWAS (89 within the window itself), and 94 were supported only by the Million Veteran Program (MVP) European ALT, AST, and NAFLD GWAS (Verma et al., 2024), which had no standalone fine-mapping, supplying the strongest colocalization for 140 genes. At the stricter prior p12 = 1 $\times$ 10$^{-6}$, 314 of 436 genes kept SuSiE PP.H4 \textgreater{} 0.5, and between the MVP and UK Biobank European GWAS of the same enzyme, 49 to 63\% of each study's colocalized gene calls recurred in the other, reusing the same liver-eQTL panel (Supplementary Note 1). Across the MVP and UK Biobank European ALT and AST pairs and the MVP and FinnGen NAFLD pair, gene-level recurrence of a colocalization call in the partner study rises with its discovery PP.H4: 0.27 (locus-cluster bootstrap 95\% interval 0.16 to 0.39; 63 tested rows) at 0.5 to 0.8, 0.52 (0.38 to 0.64; 99) at 0.8 to 0.95 and 0.67 (0.57 to 0.75; 180) at 0.95 or above. Most failures colocalize the partner GWAS with a different variant (PP.H3 \textgreater{} 0.5), which can reflect different causal variants or the different LD references used for the two studies. Genes whose best tier-1/2 PP.H4 is 0.5 to 0.8 remain provisional; recurrence alone does not qualify variant identity or establish the same mechanism. Of the 436 r2 genes, 78 had recurrent gene calls in a same-trait GWAS, 133 had a tested partner without recurrence, and 225 had no testable same-trait comparison (Supplementary Note 18). MVP was fine-mapped with an out-of-sample LD reference. SuSiE was refit to the GWAS inside each gene-by-GWAS colocalization test, in the same way for every stratum, whether or not the stratum was fine-mapped on its own (Methods). To describe where the lead colocalized variants fall, we assigned one lead COLOC variant to each of the 990 union genes (927 distinct lead variants) and classified it by overlap with transcript annotations (Figure 2C). In this gene-level summary, the lead variant lay in coding sequence or a splice site for 34 genes, in a UTR for 70, in a promoter for 104, in an intron for 522, and in intergenic sequence for 260. Of the 34 coding genes, 13 have the lead variant in their own coding sequence or splice site and 21 only in another gene's. The per-gene lead in Figure 2C need not be the variant that carries a colocalization. We therefore also checked the GWAS-side lead variant of each of the 840 r2 colocalized signal pairs against fine-mapped protein-altering or canonical-splice credible-set members of other genes (PIP $\ge$ 0.1). Four primary genes are anchored at, or in tight LD (r$^2 \ge$ 0.8) with, another gene's coding variant: C2orf16 at the GCKR variant rs1260326, CWF19L1 with an ERLIN1 variant, and RECQL4 and ZFTRAF1 with an LRRC24 variant. Their colocalization cannot be separated from that coding variant with these data. Five more genes are in moderate LD (r$^2$ 0.2 to 0.8), and 29 could not be assessed because the anchor is absent from the 1000 Genomes LD panel. Some per-gene leads are coding variants of another gene but the signal-pair anchor is not: the SAMM50 anchor has r$^2$ 0.05 with PNPLA3 rs738409, and the MAU2 anchors have r$^2$ \textless{} 0.2 with TM6SF2 rs58542926 where computable. Of the 22 direct-trait genes, C2orf16 and CWF19L1 are not separable from another gene's coding variant (Supplementary Note 18).

We describe the complete result as a map of inherited variation relevant to MASLD and liver traits, and we reserve direct MASLD language for the tiered disease subset. We restrict regulatory language to supported GWAS-eQTL relationships; protein-altering and unresolved noncoding signals keep their own labels. The eQTL side reuses an existing panel of European-ancestry donors (Broadaway et al., 2024). This analysis therefore adds MASLD and liver-trait tiering, integration of multi-ancestry GWAS with a European liver-eQTL panel, fine-mapped DNA context, and comparison with diseased human liver, not a new molecular resource. A colocalized gene is a genetically anchored candidate; colocalization alone does not show that the gene mediates the trait association.

Ancestry labels describe the GWAS populations. They do not mean that molecular regulation was ancestry-matched in every analysis. The liver-eQTL panel is European, so a non-European GWAS paired with it tests cross-population sharing, not ancestry-matched liver regulation. We report LD and eQTL provenance for these pairs and do not treat them as ancestry replications (Figures 2D and 2E). Of the 430 multi-signal genes, 334 (78\%) were supported only in European studies, 64 (15\%) in European and at least one non-European study, and 32 (7\%) only in non-European studies (23 sub-threshold and 9 genome-wide significant: 8 EAS, 1 AFR; Figure 2D). Of the 334 European-only genes, 213 were not tested in any non-European study, so their European-only status is untested rather than unsupported. Thirteen genes had PP.H4 \textgreater{} 0.5, by SuSiE or ABF, in four of the five GWAS ancestry groups, and no gene was supported in all five (Figure 2E). These partitions do not establish biological differences between ancestry groups.

\subsection{Colocalization narrows liver-trait GWAS signals to candidate genes for testing in other measurements}\label{coloc-narrows-candidates}

Together, fine-mapping and colocalization narrow the GWAS signals to a list of candidate genes that later sections test in other measurements (Figure 2A--E). Of 28,233 genes in the Gene Catalog, 18,724 could be tested for liver-eQTL colocalization, 430 colocalized by the multi-signal method after the PSD eQTL refit, and 420 of these are measured in bulk liver RNA. Colocalization marks a gene for follow-up; it does not show that the gene affects MASLD. The candidates differ in how far they can be followed (Figure 2G): 26 of the 420 change in bulk MASLD liver (Figure 3), 370 are detected in single-cell RNA, 241 in a liver proteome, 29 are on the CosMx spatial panel and 19 are detected in a plasma proteome. Most candidates take their strongest colocalization from ALT, AST or GGT rather than from a direct MASLD or liver-fat GWAS (Figure 2F). For loci colocalized with ALT or AST, the lead variant's MASLD and liver-fat associations are given per locus in Supplementary Note 22 and Figure S2Z--S2AB. Most show little detectable MASLD effect, so the Catalog reports the trait source beside each candidate and does not treat an enzyme colocalization as disease evidence.

We next asked which trait class supplies each gene's strongest colocalization evidence (Figure 2F). For 417 of the 430 genes it came from a liver-enzyme trait (AST 155, ALT 135, GGT 127), and for 13 from a direct disease or liver-fat trait (NAFLD 8, PDFF 4, NASH 1). Each gene is counted once, under the trait class of its strongest colocalization, so this split does not show how many genes each trait class supports. Counting any direct-trait support, 22 genes have a direct-trait SuSiE PP.H4 above 0.5: 12 direct only and 10 with both classes. One gene, MAU2, has PP.H4 of 1.0 in the MVP European ALT, AST and NAFLD GWAS; exact ties go to the study identifier sorting first alphabetically, so MAU2 is counted under the enzyme class (Methods). Multi-signal colocalization tested far more enzyme than direct-trait gene-by-GWAS pairs (12,121 against 771), so the predominance of liver-enzyme drivers largely follows how many gene-trait pairs were tested in each trait class (per genome-wide lead, yields are 0.134 for direct traits and 0.149 for enzymes).

\subsection{Liver-enzyme heritability, plasma hepatocyte-protein co-elevation, metabolite colocalization and TNFSF10}\label{liver-enzyme-heritability-plasma-protein-metabolite-tnfsf10}

Because 417 of the 430 primary colocalized genes take their strongest evidence from serum ALT, AST or GGT, we quantified how much serum liver-enzyme genetics is shared with MASLD, liver fat and cirrhosis (Supplementary Note 22). In a three-trait LDSC/GenomicSEM genetic covariance model (\texttt{ENZFAT1}, using sub-matrix positive-definite projection and verified against analytical delta-method standard errors), 0.701 of ALT heritability (95\% MVN interval 0.601 to 0.775; delta-method SE 0.0434; replication trio 0.802, 95\% interval 0.686 to 0.880) is shared with neither MRI liver fat nor cirrhosis (two-trait unshared fractions: 0.786 against liver fat alone and 0.769 against cirrhosis alone). Across 167 colocalized MVP European ALT/AST loci oriented to each locus's defining enzyme-raising allele (101 ALT-defined, 66 AST-defined; Figure S2Z), 19 loci carry clear MASLD liability (z $\ge$ 3.29, including AST-defined ALDH16A1/RASIP1 at z\_AST = +9.60, z\_MASLD = +4.52), 75 raise enzymes with no detectable MASLD effect (enzyme only), 73 are intermediate and 0 are opposite (overall median MASLD/enzyme z ratio 0.180, compared with 0.891 at PNPLA3 and TM6SF2). In an empirical-Bayes allele-direction analysis of 286 prioritized genes (\texttt{DIR6}, 197 independent loci; log-space half-uniform posterior evaluation passing 6 of 6 positive/negative controls with zero underflow), 17 genes (6.0\%) were concordant with liver fat or cirrhosis, 3 (1.1\%) reversed, 260 (90.9\%) had enzyme support only (most underpowered against the equivalence bound) and 6 were untestable; restricting to the non-circular European enzyme-defined subset (253 genes, 180 loci) gave 12 concordant (4.7\%), 1 reversed, 234 enzyme-only (92.5\%) and 6 untestable (with 9 LDL, 6 triglyceride, 1 T2D and 1 CAD multi-gene loci showing mixed co-effect directions at q \textless{} 0.05).

We next tested two mechanistic distinctions across the 167 enzyme loci. First, in UKB-PPP plasma proteomics (\texttt{PROT5}, 21 hepatocyte-enriched plasma proteins excluding cis loci and 21 non-hepatocyte negative-control proteins), the enzyme-raising allele elevated other hepatocyte proteins in aggregate across loci (IVW slope 0.333 protein SD per MVP ALT SD with zero UKB sample overlap, SE 0.0218, Holm \emph{p} = 5.87 $\times$ 10$^{-52}$; 0.435 per Pan-UKBB ALT SD, SE 0.0248, Holm \emph{p} = 2.45 $\times$ 10$^{-68}$; negative-control panel slope $-$0.0002 on MVP ALT, \emph{p} = 0.984), whereas individually only 26 of 167 loci (15.6\%) passed per-locus panel elevation (\texttt{panel rises}). MR-Egger confirmed a positive slope without major directional pleiotropy on non-overlapping MVP ALT (slope 0.377, \emph{p} = 1.98 $\times$ 10$^{-23}$; intercept $-$0.00137, \emph{p} = 0.065; Pan-UKBB slope 0.492, intercept $-$0.00146, \emph{p} = 0.027); the GOT1 cis control reached z = 4.77 (below the prespecified 5.45 cutoff, and no cis-GOT1 variant in UKB-PPP exceeds \textbar z\textbar{} = 5.15) while its trans panel remained null (z = $-$0.40). Second, in plasma metabolome and lipidome colocalization (\texttt{MQTL2}, 54,600 locus-trait pairs with empirical \texttt{sdY} estimation across 100\% of pairs, median 1.0002), the 6 ALT/AST loci with empirical-Bayes liver-fat support (5 steatosis-concordant plus reversed-sign PPP1R3B) colocalized with plasma lipid species (PP.H4 $\ge$ 0.8) more often than other loci (4 of 6 versus 10 of 143, odds ratio 25.2, 95\% CI 3.18 to 308.4, \emph{p} = 6.80 $\times$ 10$^{-4}$; restricted to genome-wide significant Pan-UKBB ALT loci, 4 of 5 versus 8 of 77, OR 31.7, \emph{p} = 0.0013; restricted to steatosis-concordant loci excluding PPP1R3B, 3 of 5 versus 11 of 144, OR 17.3, \emph{p} = 0.0060; attenuated to OR 5.45, \emph{p} = 0.053 at PP.H4 $\ge$ 0.5). At the worked enzyme-only locus TNFSF10 (Figure S2AB; \texttt{TRAIL4}), PSD SuSiE-coloc confirms shared causal variation between ALT and plasma TRAIL (PP.H4 = 1.000) and platelet count (PP.H4 = 1.000, distinct from albumin PP.H3 = 0.937), exact four-trait Wakefield ABF yields posterior probability 0.971 for a shared ALT--TRAIL--platelet--albumin configuration (0.943 at p12 = 5 $\times$ 10$^{-6}$; HyPrColoc 0.869), and zero variants at r$^2 \ge$ 0.8 overlap any TNFSF10 exon, coding sequence or splice site (maximum CDS r$^2$ = 0.012), ruling out a coding epitope artifact.

\subsection{The Resource assigns each gene a genetic, disease-state or two-axis evidence class}\label{the-resource-assigns-each-gene-a-genetic-disease-state-or-two-axis-evidence-class}

Having defined the colocalized gene set, we compared it with the disease-state transcriptome from the bulk fit over 774 genetic-identity clusters, which are not verified donors (Figure S1A). Among the 3,587 stable-Ensembl genes evaluated by multi-signal colocalization after shared-posterior eligibility qualification and the PSD eQTL refit, 420 were SuSiE-colocalized and 209 were canonical disease-state genes. The two sets shared 26 genes against 20.9 expected within expression deciles (95\% null range 13 to 29; Mantel-Haenszel OR 1.321, 95\% CI 0.851 to 2.051; exact \emph{p} = 0.222, one original test with updated eligibility). Historical expression diagnostics were not rerun in this qualified universe, and we state no direction. The count does not establish that the two layers are biologically separate or convergent. A further 10,928 genes in the native RNA--genetics interface, including 935 canonical disease-state genes, have no multi-signal result; 861 of those 935 lack a liver-eQTL credible set. Another 189 genes, 11 with a canonical RNA call, are untestable: 187 after the shared-posterior check and 2 after the PSD eQTL refit. Neither group is negative. Another 699 canonical disease-state genes, 503 of them lncRNAs, have no tier-1/2 colocalization record and lie outside the 14,704-gene native interface. These exclusions leave 209 of the 1,854 calls in the evaluated comparison. The raw symbol-grouped universe has 14,720 rows and the annotation-adjudicated Catalog has 14,502 jointly eligible identities; their denominators are distinct.

Across the 3,587 evaluated genes, PP.H4 and the absolute disease-state log2 fold change had Spearman rho = 0.00094 (nominal \emph{p} = 0.955; Figure S3Z). Because PP.H4 is unsigned, we did not test it against the signed fold change. Differential expression cannot serve as a required validation filter for genetically anchored regulatory candidates, and absent genetic support does not make a disease-state program unimportant: colocalization tests inherited regulation and has limited power. The Gene Catalog keeps assay applicability separate from support. Genetically anchored genes have SuSiE PP.H4 \textgreater{} 0.5 (or a Tier-1 direct-MASLD/PDFF fine-mapped coding variant) and an applicable RNA comparison without a canonical call. Established-state associated genes have a canonical RNA call and an applicable genetic comparison without genetic support. Genes supported on both axes meet both thresholds; genetic and RNA directions are not compared. Unresolved genes are testable in both domains and meet neither threshold; genes untestable in either domain are untested (Methods).

We next checked the colocalization state of the recurring disease-remodeling examples. None of GNMT, MAT1A, CYP2C19, IGFBP7, and BICC1 has a primary SuSiE PP.H4 above 0.5. Only GNMT has a single-signal ABF PP.H4 above 0.5, in one direct-NAFLD stratum (0.706, European NAFLD GWAS; Ghodsian et al., 2021). Its shared-variant anchor is at GRCh37 6:42931261 (rs2296804, SNP.PP.H4 0.259); SNP.PP.H4 is the posterior that a given variant is the shared one. The NAFLD association at this locus is weak: the anchor has GWAS p = 2.6 × 10⁻⁴, and the smallest p within 1 Mb is 1.4 × 10⁻⁴. The colocalization therefore rests mainly on the strong eQTL (p = 2.9 × 10⁻⁶⁶) and on the prior probability that one variant affects both traits (p12 = 5 × 10⁻⁶). In the earlier run of this pair (PP.H4 0.708 at that prior), PP.H4 was 0.33 at p12 = 1 × 10⁻⁶ and 0.83 at 1 × 10⁻⁵. We therefore treat GNMT's genetic support as sub-threshold sensitivity evidence. Multi-signal SuSiE colocalization was attempted for GNMT in that stratum but returned no credible-set result, and its best SuSiE PP.H4 in any primary study is 0.128 (UK Biobank ALT). The five large European enzyme GWAS give GNMT ABF PP.H3 of 0.94 to 1.00, but in four of them this reflects a separate association about 0.8 Mb away (r² \textless{} 0.005 with the NAFLD anchor); conditioning the NAFLD association on each of these distant leads leaves PP.H4 at 0.64 to 0.69. In a post hoc analysis restricted to 250 kb around the anchor, those four GWAS have no association strong enough to call a shared or a distinct variant. UK Biobank ALT has a local association (p = 1.5 × 10⁻¹³, 14 kb from the anchor) that does not colocalize with the GNMT eQTL (PP.H3 0.88), and MVP European ALT shows no association at the anchor (p = 0.70) (Supplementary Note 18). The strongest stored ABF anchors for MAT1A, CYP2C19, IGFBP7, and BICC1 are at 10:81143387, 10:96292988, 4:57949517, and 10:60583832, respectively, but their gene-trait PP.H4 values remain below 0.5. These coordinates are COLOC shared-variant anchors, not causal variants or GWAS credible-set PIP.

We treat GNMT as a provisional worked locus. At the nonpalindromic variant rs2296805, the G allele lowers liver GNMT expression and aligns with higher direct-NAFLD risk. rs2296805 is a proxy for the rs2296804 anchor (UK Biobank European r2 = 0.998). The anchor is a C/G variant at a frequency near 0.5, so its allele labels do not resolve strand on their own. The nonpalindromic rs9471976 also aligns lower expression with higher risk and has r2 = 0.818 with rs2296805. rs11752813 (r2 = 0.580) is palindromic, and its reported sign is not independently harmonized. Ji et al.~(2012) associated rs9471976 and rs11752813 with GNMT protein in 268 liver biopsies and reported allele-dependent reporter activity, which motivates resolving this locus. Unsigned LD and colocalization do not identify a causal variant or show that GNMT mediates the association.

Expression-based colocalization may miss some protein-altering and protein-level loci. Protein-altering variants can act through protein function without changing steady-state RNA, so they need not appear as expression QTLs; others change RNA abundance through nonsense-mediated decay or splicing (Rivas et al., 2015). SuSiE credible sets place a missense variant at three such genes: MTARC1 (rs2642438; normalized posterior inclusion probability, PIP, 0.757 in a European PDFF GWAS), TM6SF2 (rs58542926; PIP 1.0, European PDFF), and PNPLA3 (rs738409; PIP 0.706, MVP Admixed American NAFLD). External protein-QTL evidence for MTARC1 is a source finding (Gobeil et al., 2026b). In total, 31 genes carry a fine-mapped protein-altering or canonical-splice credible-set member at normalized SuSiE PIP $\ge$ 0.5. Restricting coding-driven promotion to Tier-1 direct MASLD/PDFF traits retains 9 genes (GCKR and SAMM50, already supported by regulatory colocalization, plus 7 newly promoted genes: ADH1B, APOE, MTARC1, PNPLA2, PNPLA3, SERPINA1 and TM6SF2), moving those 7 unresolved genes to genetically anchored while keeping all 17 non-regulatory Tier-2 liver-enzyme-only coding hits (including HBB, GOT1, ABCB11, EPPK1 and FOXN4) in \texttt{coding\_liver\_enzyme\_gwas\_only} outside \texttt{genetically\_anchored} and \texttt{both\_supported} (Figures 6B and 6C). This field misses non-canonical splice effects and does not show that the gene mediates the trait association.

\subsection{Noncoding variants carry most fine-mapped posterior probability at liver-trait loci}\label{noncoding-variants-carry-most-fine-mapped-posterior-probability-at-liver-trait-loci}

To measure DNA architecture directly, we used every reliable tier-1/2 SuSiE credible set that fine-mapping produced, not only the credible sets that can be rebuilt from the top SNP of a colocalized gene. Of 6,467 study-by-locus credible-set instances, 6,415 passed the prespecified PIP-sum criterion. Figure S2X uses the 4,107 of these whose lead variant has p \textless{} 5 × 10⁻⁸ in its own GWAS (164 direct-trait and 3,943 enzyme credible sets; Methods). Complete GRCh37 reference-oriented VEP annotation (McLaren et al., 2016) assigned 100\% of their normalized PIP to a consequence class. We grouped credible sets into physical regions by single linkage, joining sets whose lead positions lie within 1 Mb of each other, so one region can span more than 1 Mb, and weighted each region equally. Direct MASLD and liver-fat loci then allocated 29.42\% of the probability to protein-altering variants, 0.23\% to canonical splice variants, 3.92\% to synonymous or UTR variants, and 66.43\% to other noncoding variants. Liver-enzyme loci allocated 6.67\%, 0.00\%, 3.50\%, and 89.83\% (Figure S2X). The grouping limits the weight of repeatedly recovered studies and does not assert LD or statistical independence. When each credible set was instead weighted equally, protein-altering variants carried 17.36\% of the posterior for direct traits and 3.43\% for liver enzymes.

Fine-mapping coverage differs by stratum. Every lead window of the European direct-trait GWAS that entered fine-mapping has a credible set (47 of 47). The European enzyme arm consists of 82 credible sets from the UK Biobank ALT, AST, and GGT GWAS (Sinnott-Armstrong et al., 2021). In the fine-mapping table used for Figure S2X, these GWAS have credible sets at 6, 9, and 15 loci and none at PNPLA3 or TM6SF2. These loci came from an earlier lead set, so they match 8 of the 767 current lead windows by locus identifier, while 23 of 767 have a credible set within 500 kb by position, the better coverage measure. A separate UKBB fine-mapping run at loci selected by colocalization (PP.H4 ≥ 0.5) is not used, because loci selected on colocalization cannot serve as an unbiased enzyme arm. The MVP European ALT, AST, and NAFLD GWAS had no standalone fine-mapping. In the Admixed American (AMR) GWAS, all 4 direct and all 35 enzyme lead windows at p \textless{} 5 × 10⁻⁸ have a credible set.

Taking protein-altering and canonical-splice variants together, direct and enzyme loci pooled over ancestries carry 29.65\% and 6.67\% of their posterior probability, respectively (difference 0.2298, 95\% physical-cluster bootstrap interval 0.0808 to 0.3958; 120 physical clusters: 16 direct, 115 enzyme, and 11 shared). For this interval, we resampled whole clusters 10,000 times, because nearby credible sets are not independent. This pooled value is not a trait-class contrast. Of the direct credible sets, 106 of 164 are European, and 15 of the 16 direct clusters contain a European direct set, whereas 3,861 of 3,943 enzyme credible sets are non-European. Trait class is therefore confounded with GWAS ancestry, LD reference, and fine-mapping coverage.

To separate trait class from ancestry, we stratified the comparison by ancestry (Figure S2Y; 10,000 shared draws over 44 pooled clusters). European studies gave 29.88\% against 6.92\% (difference 0.2296, 95\% interval 0.0563 to 0.4211), and AMR studies gave 8.11\% against 6.91\% (0.0120, −0.1246 to 0.1620). The prespecified fixed-weight European-plus-AMR estimate (weights 16/20 and 4/20, the direct-cluster counts) was 0.1861 (0.0503 to 0.3337). All non-European studies pooled gave 8.71\% against 6.48\% (0.0223, −0.0831 to 0.1686), a descriptive value with no calibration control. None of these estimates receives a trait-class call. The European estimate failed two checks: coverage (enzyme credible sets at 8 of 767 lead windows by locus identifier and 23 by position, direct arm at 47 of 47) and power (minimum detectable difference 0.303 against the pooled 0.230). It passed its random-split calibration control. In this control, we split East Asian enzyme clusters (East Asian and African clusters for the European-plus-AMR estimate) at random into two pseudo-classes with the observed cluster counts, 200 times. Both pseudo-classes come from one trait class, so a calibrated 95\% interval should cover zero in at least 180 splits, and the 98.333\% interval (Bonferroni over the three estimates) in at least 190. The European 95\% and 98.333\% intervals covered zero in 183 and 191 splits. AMR coverage passed (4 of 4 direct and 35 of 35 enzyme lead windows), but the AMR and European-plus-AMR estimates failed the control (122 and 172 splits at 95\%; 135 and 188 at 98.333\%), so their intervals are too narrow. Both also narrowly failed the power check (minimum detectable differences 0.234 and 0.232), and in 156 of the 10,000 draws the AMR direct arm, 4 clusters, was empty, more than the prespecified limit of 100. The European-plus-AMR estimate also carries the European coverage failure.

Within the direct class, the three European PDFF GWAS (37 credible sets, 12 clusters) carry 40.14\% against 6.67\% for all enzyme traits (difference 0.3347, 95\% interval 0.1492 to 0.5280). This comparison sets European PDFF studies against an enzyme arm that is mostly non-European, so it carries the same ancestry confounding as the pooled value. NAFLD plus NASH (127 sets, 8 clusters) carry 15.92\% (0.0925, −0.0596 to 0.2844) (Figure S2S). The pooled direct-minus-enzyme difference is also concentrated. Five of the 16 direct clusters carry 73.9\% of the summed direct cluster fractions, and the largest, one PDFF credible set at chr11:823,586, carries 21.1\%. The 50 robustness comparisons (study and direct-cluster omissions, five grouping windows, and one LD-block comparison) gave ratios from 3.44 to 5.18 around the pooled reference ratio of 4.45, a descriptive value. These comparisons include no ancestry or trait stratification and are correlated sensitivities, not replications.

The two trait classes in the 29.65\% against 6.67\% comparison also differ in GWAS size and effect size. Direct-trait credible sets come from smaller GWAS than enzyme credible sets (median effective n 21,893 against 160,000 over the 4,134 credible sets at genome-wide leads) and carry larger standardized lead-variant effects (median 0.0423 against 0.0190). Supplementary Note 4 reports a sensitivity analysis matched on standardized lead-variant effect, over the same 4,134 credible sets, including those that failed the PIP-sum criterion. No permutation \emph{p} value is reported. Comparing the trait classes within one ancestry requires SuSiE fine-mapping of all 767 UK Biobank ALT, AST, and GGT lead windows and of the MVP European ALT, AST, and NAFLD GWAS. This comparison does not test whether the pooled coding difference contributes to the predominance of liver-enzyme evidence among colocalized genes, or to which genes appear in both the colocalized and disease-state sets.

Within the noncoding posterior probability, we treated promoter, liver ABC-enhancer (Nasser et al., 2021), and liver-lineage open-chromatin annotations as overlapping contexts, not as probabilities that must sum to one. After re-annotation by genomic position (Supplementary Note 6), no noncoding posterior probability was left without a context call. Summed over all credible sets of each class at genome-wide leads, 14.93\% of direct-trait and 14.01\% of liver-enzyme noncoding posterior probability lies within 2 kb of a TSS, 10.73\% and 7.84\% in liver-lineage open chromatin, 0.30\% and 3.85\% in a liver ABC enhancer, and 75.39\% and 78.28\% in sequence with none of these contexts. These pooled values mix ancestries. AMR is the one ancestry with comparable fine-mapping coverage in both classes. Within AMR, no context difference received a call (99\% intervals, Bonferroni over five contexts). Every random-split calibration control failed, and in 167 of the 10,000 draws the AMR direct arm, 4 clusters, was empty, more than the prespecified limit of 100. For the share with no deposited context, direct minus enzyme was −0.209 (99\% interval −0.599 to 0.166). Thus, most noncoding fine-mapped probability lies outside the deposited promoter, open-chromatin, and ABC annotations, and the DNA result does not assign a regulatory mechanism to most noncoding signals.

\subsection{A bulk-liver accessibility QTL shares the GWAS signal at 191 of 581 colocalized lead signals}\label{a-bulk-liver-accessibility-qtl-shares-the-gwas-signal-at-191-of-581-colocalized-lead-signals}

We next asked whether some of these signals share a variant with an accessible chromatin region in bulk liver. We took 581 European tier-1/2 GWAS lead signals that colocalize with a liver eQTL. For each, we tested whether a chromatin accessibility QTL measured in bulk liver from 138 donors (Currin et al., 2025) shares the same GWAS signal (PP.H4 \textgreater{} 0.8 and LD r² ≥ 0.5 between the two lead variants). Of these signals, 191 share a signal with a bulk-liver accessible region, and Currin et al.~(2025) already reported at least 45 of them for the same trait and accessible region. Another 315 map to a different variant, 60 had no testable accessible region, 9 colocalized without confirmed lead LD, and 6 were indeterminate. Of the 436 primary SuSiE genes, 143 have a colocalized accessible region. In 108 of them, the same accessible region also colocalizes with the gene's eQTL. This comparison uses no lead-variant LD rule. When the accessibility-QTL and eQTL lead variants must also have r² ≥ 0.8, 100 of the 108 remain (107 at r² ≥ 0.5). The other 35 genes (43 under the r² ≥ 0.8 rule) are reached only through the GWAS-eQTL pair. Two pairwise colocalizations with one GWAS signal do not show that the accessibility QTL and the eQTL share a variant, so we do not link these genes to the region. The eQTL and accessibility-QTL donor sets may overlap, so the three-way support is not independent replication. At 88 three-way loci with a nonpalindromic variant, the allele that raises accessibility also raises the gene's liver expression at 67 (0.76; binomial Holm p = 1.8 × 10⁻⁶ over two tests), as expected if the accessible region acts as an enhancer. Whether the trait-raising allele moves expression in the disease-state direction was not shown (28 of 49 loci, Holm p = 0.39) (Supplementary Note 18).

To estimate how often a shared accessible region occurs, we excluded chromosomes 1 and 19, which a pilot analysis had examined. Of the remaining 238 independent liver-trait loci, 80 (0.336, Wilson 95\% CI 0.279 to 0.398) carry a shared accessible region. This rate is conditional on prior liver-eQTL colocalization and is not the share of liver GWAS loci with an accessibility QTL. We then compared loci selected the same way that carry signals from one trait arm only (172 liver-trait loci and 78 platelet-count loci). At equal GWAS and accessibility-QTL strength, liver-trait signals shared a liver accessibility QTL more often than platelet-count signals (adjusted odds ratio 3.35, 95\% CI 1.62 to 7.40; \emph{p} = 0.0016, one prespecified test). The unadjusted counts, 70 of 172 against 11 of 78, are descriptive. This agrees with the trait-level pattern reported by Currin et al.~(2025), and we do not present it as new.

We used one sequence model, AlphaGenome 2 (evaluated using the \texttt{ALL\_FOLDS} checkpoint, whose pretraining may overlap liver chromatin tracks), only to annotate predicted accessibility at credible-set variants, after a benchmark with endpoints fixed in advance. Without fitting, its predicted bulk-liver ATAC effects agreed with measured liver caQTL effects at 3,845 lead variants (signed Spearman 0.719, 95\% CI 0.704 to 0.733; version 1 on the same leads, 0.736; hepatocyte single-nucleus ATAC, 0.747; Supplementary Note 21). Its expression readouts did not predict liver eQTL direction (Supplementary Note 21), so the model supplies no predicted gene and no predicted expression direction. Across 4,639 credible-set members, applying the 99th-percentile matched-background accessibility rule flags 126 members across three tracks (2.17\% background rate, ~101 expected; locus-grouped PIP permutation confirms modest enrichment, B1a Holm \emph{p} = 0.018 and B1b Holm \emph{p} = 0.050, while BH over per-variant scores retains 4 of 24 track calls at q \textless{} 0.05; in measured Currin bulk caQTL overlap [n = 80 measured variants], 26 of 26 sequence-supported variants versus 30 of 54 others overlap, exact \emph{p} = 1.23 $\times$ 10$^{-5}$, Mantel-Haenszel Holm \emph{p} = 0.0051). The label is therefore an accessibility annotation shown beside the measured colocalization, caQTL and fine-mapping evidence. For any single variant it is weak evidence, and it does not identify a causal variant.

Taken together, these analyses define the genetic arm of the Resource as a bounded map of trait-associated inherited variation, fine-mapped molecular context, and colocalized liver-trait regulatory candidates. This map does not substitute for the established disease-state transcriptome, which we describe next.
